\section{Contents of this supplementary material}\label{app:overview}

Section letters below are this document's; numbered sections, equations, tables and figures
cited without a letter are the main paper's. One naming difference: this supplementary keeps
the long-form names for the two parts of a gain. The main paper's \emph{group-level} part
$\dP$ is called the \emph{transported gain} here, and its \emph{case-level} part $\dR$ the
\emph{within-group covariance gain}; the quantities are identical. \Cref{app:formal} states the eleven formal results
the paper uses and \cref{app:proofs} proves them. The remaining sections give the worked
four-point example (\cref{app:worked}), the influence-function derivation
(\cref{app:influence}), the randomization algorithm (\cref{app:random}), implementation,
architectures, seeds and compute (\cref{app:implementation}), the sensitivity bound evaluated
on our own estimates (\cref{app:sensitivity}), the conditions of the limit theorem
(\cref{app:clt}), the choice of grouping variable (\cref{app:partition}), a terminology table
(\cref{app:terms}), the four-point example showing that the group-level part is not prior
fit (\cref{app:transported}), and the full protocols and tables for the controlled
predictors, the CLIP comparison, the sensitivity analyses and the two long-tailed studies
(\cref{app:studies}), and the mean-recall split of the released checkpoint with the
evaluator replay, the operating-point control, the within-cell AUC and the proper-score split
in every configuration (\cref{app:recall}), the CLIP study over five seeds (\cref{app:vgseeds}),
and the scene-graph audit under 20 calibration halves, merged predicates, a label shift and a
zero-shot CLIP model (\cref{app:robust}).

The code, the results files behind every number, the scripts that regenerate each model's
predictions, and a script that checks every number printed in the paper against the results
files are uploaded as an anonymized archive alongside this document.
